
Logistic regression trained on six GitHub metadata features achieves 95.8--100\% accuracy for predicting maintenance status of Papers with Code benchmark repositories, establishing a simplicity baseline for this domain. Using a dataset of 120 real machine learning benchmark repositories, binary classification (maintained vs. abandoned, defined by 180-day commit threshold) was performed with features extracted from the GitHub API: stars, forks, contributors, commits, issues, and days since last commit. Logistic regression with balanced class weights achieved 100\% accuracy on a held-out test set when trained with eight features (including two derived features later identified as tautological), and 95.8\% accuracy when restricted to six independently measured GitHub metadata features. Gradient boosting achieved perfect 100\% accuracy on the same six-feature dataset. Coefficient analysis revealed a two-tier signal hierarchy: staleness (days\_since\_last\_commit, coefficient -3.05) dominated the signal, approximately five times stronger than engagement metrics (forks, issues, contributors: +0.14 to +0.55). Feature importance divergence between logistic regression (multi-feature strategy) and gradient boosting (single-feature dominance: days\_since\_last with importance 1.0) indicated approximate but not perfect linear separability. The 4.2 percentage point gap between logistic regression and gradient boosting fell within the pre-specified 5\% threshold for mechanism validation. Limitations include domain specificity (Papers with Code repositories only), small sample size (120 repositories, 24-sample test set), temporal stability untested (IID evaluation only), and tautological feature contamination in initial experiments (subsequently corrected). These findings demonstrate that six GitHub metadata features are sufficient for near-perfect maintenance classification in the benchmark repository domain, without graph-based features or ensemble methods. Future work should validate generalization to non-ML repositories and assess temporal robustness.
